\chapter{微分中值定理}
\section{拉格朗日定理和函数单调性}
\begin{theorem}[费马定理]
    设函数$f$在点$x_0$处取得极值,且$f$在点$x_0$处可导,则$f'(x_0)=0$.
\end{theorem}
\begin{theorem}[罗尔定理]
    设函数$f$在闭区间$[a,b]$上连续,在开区间$(a,b)$上可导,且$f(a)=f(b)$,则存在$\xi\in(a,b)$,使得$f'(\xi)=0$.
\end{theorem}
\begin{proof}
    由最大最小值定理,$f$在$[a,b]$上存在最大值和最小值,分别设为$M$和$m$.若$M=m$,则$f(x)\equiv M$,于是任意的$\xi\in(a,b)$,都有$f'(\xi)=0$;若$m<M$,则$m$和$M$至少有一个不等于$f(a),f(b)$,于是存在$\xi\in(a,b)$,使得$f(\xi)$等于$M$或$m$,且作为极值点,根据费马定理,$f'(\xi)=0$.
\end{proof}
\begin{theorem}[拉格朗日定理]
    设函数$f$在闭区间$[a,b]$上连续,在开区间$(a,b)$上可导,则存在$\xi\in(a,b)$,使得
    \begin{equation*}
        f'(\xi)=\dfrac{f(b)-f(a)}{b-a}.
    \end{equation*}
\end{theorem}
\begin{proof}
    设函数$F(x)=f(x)-\dfrac{f(b)-f(a)}{b-a}x$,显然$F$在$[a,b]$连续,$(a,b)$可导.并且$F(b)-F(a)=f(b)-f(a)-\dfrac{f(b)-f(a)}{b-a}(b-a)=0$,即$F(a)=F(b)$.根据罗尔定理,存在$\xi\in(a,b)$,使得$F'(\xi)=0$,即$f'(\xi)=\dfrac{f(b)-f(a)}{b-a}$.
\end{proof}
\begin{theorem}[导数极限定理]
    设函数$f$在$U(x_0)$连续,在$U^{\circ}(x_0)$上可导,且$\displaystyle\lim_{x\to x_0}f'(x)$存在,则\[f'(x_0)=\displaystyle\lim_{x\to x_0}f'(x).\]
\end{theorem}
注意区分:$f(x)$在$x_0$处的导数$f'(x_0)$与导函数$f'(x)$在$x_0$处的极限$\displaystyle\lim_{x\to x_0}f'(x)$.

上述定理可以分为两部分:$f$在$U_{+}(x_0)$连续,在$U^{\circ}_{+}(x_0)$上可导,且$\displaystyle\lim_{x\to x_0^{+}}f'(x)$存在,则$f'_{+}(x_0)=\displaystyle\lim_{x\to x_0^{+}}f'(x)$;\\$f$在$U_{-}(x_0)$连续,在$U^{\circ}_{-}(x_0)$上可导,且$\displaystyle\lim_{x\to x_0^{-}}f'(x)$存在,则$f'_{-}(x_0)=\displaystyle\lim_{x\to x_0^{-}}f'(x)$.
下面证明上述定理的第一部分,第二部分的证明类似.
\begin{proof}
    任取$x\in[x_0,x]$,根据拉格朗日定理,存在$\xi\in(x_0,x)$,使得\[f'(\xi)=\dfrac{f(x)-f(x_0)}{x-x_0}.\]当$x\to x_0^{+}$时,$\xi\to x_0^{+}$,所以$\displaystyle\lim_{x\to x_0^{+}}f'(\xi)=\displaystyle\lim_{\xi\to x_0^{+}}f'(\xi)=\displaystyle\lim_{x\to x_0^{+}}f'(x)$,而$\displaystyle\lim_{x\to x_0^{+}}\dfrac{f(x)-f(x_0)}{x-x_0}=f'_{+}(x_0)$.于是有$\displaystyle\lim_{x\to x_0^{+}}f'(x)=f'_{+}(x_0)$.
\end{proof}
\begin{theorem}
    设函数$f$在$(a,b)$可导.证明:$\forall x_0\in(a,b)$,$x_0$不是$f'(x)$的第一类间断点.(导数没有第一类间断点)
\end{theorem}
\begin{proof}
    由$f$在$(a,b)$可导,可知$f$在$(a,x_0]$与$[x_0,b)$上连续,$f$在$(a,x_0)$与$(x_0,b)$可导.若$x_0$是$f'(x)$的第一类间断点,则$\displaystyle\lim_{x\to x_0^{+}}f'(x),\displaystyle\lim_{x\to x_0^{-}}f'(x)$均存在,根据导数极限定理,有$f'_{+}(x_0)=\displaystyle\lim_{x\to x_0^{+}}f'(x),f'_{-}(x_0)=\displaystyle\lim_{x\to x_0^{-}}f'(x)$.\\又$f$在$x_0$可导,所以$f'(x_0)=f'_{+}(x_0)=f'_{-}(x_0)$,从而有$f'(x_0)=\displaystyle\lim_{x\to x_0^{+}}f'(x)=\displaystyle\lim_{x\to x_0^{-}}f'(x)$,说明$f$在$x_0$处连续,与$x_0$是$f'(x)$的第一类间断点矛盾.因此,$x_0$不是$f'(x)$的第一类间断点.
\end{proof}